\documentclass[11pt]{article}
\usepackage[margin=25mm]{geometry}
\usepackage{fontspec}
\setmainfont{Latin Modern Roman}
\usepackage{amsmath,amssymb,mathtools}
\usepackage{unicode-math}
\setmathfont{Latin Modern Math}
\usepackage{xurl,hyperref}
\hypersetup{colorlinks=true,urlcolor=blue}
\setcounter{secnumdepth}{0}
\setlength{\parindent}{0pt}
\setlength{\parskip}{5pt}
\setlength{\emergencystretch}{3em}
\newcommand{\sref}[2]{\hyperlink{source:#1:#2}{[S#2]}}
\newcommand{\cataloguescope}[1]{\paragraph{Recorded target.}#1}
\newcommand{\R}{\mathbb R}
\newcommand{\C}{\mathbb C}
\newcommand{\N}{\mathbb N}
\newcommand{\Z}{\mathbb Z}
\newcommand{\Q}{\mathbb Q}
\newcommand{\F}{\mathbb F}
\newcommand{\eps}{\varepsilon}
\begin{document}
\section{Riemann-Hilbert Problem with non-analytic data}\label{statement}

\subsection{Definitions and mathematical statement}

\paragraph{Shared notation.}$\mathbb N=\{1,2,\ldots\}$; the standard number fields have their usual meanings. All additional conventions are stated locally.


\paragraph{Definitions and notation.}A matrix Riemann--Hilbert problem asks for a matrix function analytic off an oriented contour, with prescribed jumps $M_+(z)=M_-(z)J_n(z)$ and normalization. The oscillatory data contain $e^{in\phi(z)}r(z)$, with $n\to\infty$. Nonlinear steepest descent deforms the contour and analyzes local and global model problems. The source asks to relax analyticity of the phase, both in equilibrium-measure analysis and contour deformation.

\paragraph{Catalogue wording retained for traceability.}
In many situations one is concerned with the asymptotic behavior of Riemann-Hilbert problems with exponentially varying data of the form $e^{in\phi(z)} r(z)$, $n \to \infty$. The Deift-Zhou nonlinear steepest descent method for such problems requires $\phi(z)$ to be analytic. The analyticity is used in two ways: to control the equilibrium measure associated with the problem, and then to deform the contour for the RHP. By contrast, the method only requires minimal smoothness for $r(z)$ (see eg. Deift-Zhou in the context of Problem 12 below). It is of considerable theoretical and practical interest to extend the nonlinear steepest descent method to situations where $\phi$ is no longer analytic, and has, for example, only a finite number of derivatives. For very interesting work on the analyticity problem, we refer the reader to a recent paper of Miller and McLaughlin. There is also interesting, older work due to Varzugin.
\paragraph{Scope and source qualification.}This method-development problem leaves the minimal differentiability and error estimates to be determined. The source mentions Miller--McLaughlin and Varzugin work as partial context, not as a general solution.
\subsection{Short English statement}
Extend nonlinear steepest descent for oscillatory Riemann--Hilbert problems to phases with only finitely many derivatives.
\subsection{Research attempt}
\paragraph{Model, date and investigated scope.}
OpenAI GPT-6 Astra (configured model \texttt{gpt-6-astra}), ultra reasoning; generated 11 October 2026. Derived an explicit almost-analytic extension error and n\textasciicircum{}\{-k\} transverse integral bound on nonstationary patches. Equilibrium measures, critical points and matrix inverse estimates remain unresolved.

\paragraph{Sources and literature check.}
Read the supplied complete definition, catalogue wording, scope qualifications and source list for this canonical ID. The original-source citations are inherited as supplied references; no fresh primary-source or novelty search was performed for this attempt. No current open-status claim is inferred from this calculation.

\paragraph{Finite-smoothness extension of a phase.}Let a real phase \(\phi\) be \(C^{k+1}\) on a compact interval, where \(k\ge1\), and suppose \(\phi'(x)\ge c>0\). For \(z=x+iy\) define the truncated extension
\[\Phi(x,y)=\sum_{j=0}^{k}\frac{(iy)^j}{j!}\phi^{(j)}(x),\qquad \bar\partial=\tfrac12(\partial_x+i\partial_y).
\]
Differentiating term by term cancels all adjacent terms and leaves
\[\bar\partial\Phi=\frac{(iy)^k}{2k!}\phi^{(k+1)}(x).
\]
For sufficiently small \(y\ge0\), uniformly on a smaller compact interval, \(\operatorname{Im}\Phi\ge cy/2\): the linear imaginary term is \(y\phi'\), and all remaining terms are bounded by a constant times \(y^2\) (or vanish for \(k=1\)). Hence
\[|\bar\partial e^{in\Phi}|\le C n y^k e^{-ncy/2},\qquad
\int_0^{y_0}|\bar\partial e^{in\Phi}|dy\le C' n^{-k}.
\]
The last estimate follows by substituting \(s=ny\); the integral of \(s^ke^{-cs/2}\) is finite.
\paragraph{Why the local estimate is insufficient.}This gives a concrete replacement for exact analyticity on a nonstationary patch. It does not bound the inverse of a matrix singular-integral problem. In particular the Cauchy kernel \((z-w)^{-1}\) is singular, so a small area \(L^1\) norm of a source does not by itself imply a small \(L^\infty\) Cauchy transform. A source concentrated in a shrinking disk already warns against that inference. Stationary points, where \(\phi'=0\), also remove the uniform exponential factor used above.
\paragraph{First gap.}A full method must construct suitable extensions near stationary points and equilibrium-measure endpoints, then prove an operator-norm estimate for the resulting matrix \(\bar\partial\) problem and its interaction with local models. Neither the equilibrium-measure regularity nor those uniform inverse bounds are obtained here. The calculation is an elementary almost-analytic extension estimate, not a general finite-regularity steepest-descent theorem.

\paragraph{Verification and adversarial audit.}
Checked cancellation in dbar, differentiability order, positivity of the imaginary part, n scaling and failure of L1 control to bound a singular Cauchy transform. This was an LLM adversarial audit, not independent human verification; no theorem prover was used. Prompts A then B applied to this ID: Checked cancellation in dbar, differentiability order, positivity of the imaginary part, n scaling and failure of L1 control to bound a singular Cauchy transform. Scope remains limited: Control the matrix dbar integral operator uniformly, including critical points and equilibrium-measure endpoints.

Coordinator checked adjacent cancellation in the truncated phase extension, the imaginary-part bound and rescaled integral estimate. Singular Cauchy-operator control is left separate. This additional check was model-assisted, not independent human verification.

\paragraph{Tools, process and contributions.}
Python 3 and local shell used to read supplied files and serialize this manuscript; mathematics derived symbolically in the text, with no numerical experiment or formal proof assistant. Contribution prompt A was followed by prompt B. GPT-6 Astra with ultra reasoning drafted the research and performed its stated self-audit; any additional coordinator contribution is identified in the verification and audit. A GPT-6-based Codex coordinator assembled the manuscript and checked repository integration; its exact variant and reasoning configuration were not exposed. No human mathematical contribution, independent human proof check or formal proof certification is supplied. Complete derivations and the consulted source references are included here; private reasoning transcripts are not published.

\paragraph{Final status and first remaining gap.}
\textbf{UNRESOLVED}. Control the matrix dbar integral operator uniformly, including critical points and equilibrium-measure endpoints.

\section{Completion Estimate}
\noindent\textbf{COMPLETION: 0\%}
The calculation is an elementary benchmark or conditional reduction; it establishes no claimed new portion of the complete source target.

\subsection{Sources}
\begin{itemize}
\item[\hypertarget{source:statement:1}{S1}] ULAM.AI and the UnsolvedMath Contributors, UnsolvedMath, catalogue entry. Catalogue curation under CC BY 4.0; definitions and English exposition added locally.
\url{https://huggingface.co/datasets/ulamai/UnsolvedMath}.
\item[\hypertarget{source:statement:2}{S2}] Percy Deift, Some open problems in random matrix theory and the theory of integrable systems (2007), arXiv:0712.0849, Problem 3.
\url{https://arxiv.org/abs/0712.0849}.
\end{itemize}
\label{end:statement}
\end{document}
